\section{Smooth limits from lattice differences}
\label{ex-smooth-limits}

The extraction argument initially gives continuity only on generic
configuration patches. A stronger family of difference bounds removes
coordinate ties and identifies every derivative of the limit. This
argument also explains the additional input needed for convergence along
the whole mesh sequence: pointwise Cauchy convergence on a dense set,
without another subsequence selection.

We state the analytic result on an arbitrary open set
\(U\subset\mathbb R^d\). In the fixed-period application,
\(d=4n\), \(U\) is the set of configurations with no coincident
points modulo the period, and the functions are periodic lifts of lattice
correlations. Thus \(U\) includes configurations with coordinate ties
and can cross a fundamental-domain seam. Write
\(\Gamma_N=a_N\mathbb Z^d\), where \(a_N>0\) and \(a_N\to0\).
For functions \(F_N:\Gamma_N\to\mathbb C\), define
\[
 \Delta_{N,j}F_N(x)=\frac{F_N(x+a_Ne_j)-F_N(x)}{a_N},
 \qquad
 F_{N,\alpha}=\Delta_{N,\alpha_1}\cdots
                    \Delta_{N,\alpha_k}F_N
\]
for a finite list \(\alpha\) of coordinates, including the empty list.
The order convention is that prefixing a coordinate applies one further
difference to the previous function. Choose cell representatives
\(Q_N:\mathbb R^d\to\Gamma_N\) by rounding each coordinate, so
\(|Q_Nx-x|_1\le da_N\) and \(Q_Nx=x\) on \(\Gamma_N\).

\begin{theorem}[Full-sequence smooth convergence]\label{ex-smooth-grid-limit}
Assume that for every compact \(K\subset U\) and every finite list
\(\alpha\), there are \(C_{K,\alpha}<\infty\) and \(N_{K,\alpha}\)
such that
\begin{equation}\label{ex-grid-derivative-budget}
 |F_{N,\alpha}(x)|\le C_{K,\alpha}
 \quad(N\ge N_{K,\alpha},\ x\in\Gamma_N\cap K).
\end{equation}
Suppose also that there is a dense subset \(D\subset U\) for which
\(F_N(Q_Nx)\) is Cauchy for every \(x\in D\).
Then there is a smooth function \(f:U\to\mathbb C\) such that, for
every \(\alpha\),
\begin{equation}\label{ex-grid-all-limits}
 F_{N,\alpha}(Q_Nx)\longrightarrow\partial^\alpha f(x)
 \quad\text{uniformly on every compact subset of }U.
\end{equation}
No further subsequence is taken. In particular, whenever \(N_k\to\infty\)
and \(x_k\in\Gamma_{N_k}\) tends to \(x\in U\),
\(F_{N_k,\alpha}(x_k)\to\partial^\alpha f(x)\).
\end{theorem}
\begin{proof}
We give the details in the \(\ell^1\) norm; equivalence of norms in
fixed dimension gives the same local convergence in the Euclidean norm.
All estimates below are made in a compact neighborhood strictly inside
\(U\), and finitely many thresholds in
\eqref{ex-grid-derivative-budget} can be replaced by their maximum.

\emph{Grid Lipschitz and Taylor estimates.}
Fix \(\alpha\), and abbreviate \(G_N=F_{N,\alpha}\).
On a slightly enlarged compact set, bound all its first differences
by \(B_1\) and all its second differences by \(B_2\). Join sufficiently
nearby grid points \(x,y\) by moving one coordinate at a time, always
monotonically toward \(y\). The path has total length
\(R=|y-x|_1\) and stays in the coordinate box between \(x,y\).
Each edge changes \(G_N\) by at most \(a_NB_1\), including a
negative edge by using the forward difference at its lower endpoint.
Telescoping gives
\begin{equation}\label{ex-grid-lipschitz}
 |G_N(y)-G_N(x)|\le B_1|y-x|_1.
\end{equation}
Apply the same argument to each \(\Delta_{N,j}G_N\). At every base
point \(z\) of an edge of the path,
\(|\Delta_{N,j}G_N(z)-\Delta_{N,j}G_N(x)|\le B_2R\).
Subtract the linear function whose coordinate coefficients are
\(\Delta_{N,j}G_N(x)\) from the edge increments, and telescope again:
\begin{equation}\label{ex-grid-taylor}
 \left|G_N(y)-G_N(x)-\sum_j(y_j-x_j)\Delta_{N,j}G_N(x)\right|
 \le B_2|y-x|_1^2.
\end{equation}
The slightly enlarged compact set ensures that all shifted points used
by these differences remain inside the region of the bounds. This is
possible because only finitely many shifts occur and \(a_N\to0\).

\emph{Extension from dense pointwise convergence.}
Equation~\eqref{ex-grid-lipschitz} and the rounding bound imply on any
compact patch, for sufficiently nearby \(x,y\),
\[
 |G_N(Q_Nx)-G_N(Q_Ny)|\le B_1(|x-y|_1+2da_N).
\]
Thus the rounded functions are asymptotically equicontinuous; they need
not themselves be continuous. For \(\alpha=\varnothing\), this and
the Cauchy property on \(D\) imply uniform Cauchy convergence on each
compact set. To see it explicitly, cover that compact set by finitely
many sufficiently small balls with centers in \(D\), lying in a common
compact neighborhood in \(U\). At the centers, choose one common
Cauchy threshold. Choose the mesh still finer so that the two oscillation
errors, from a point to its center for each of the two mesh functions,
are small. The triangle inequality then bounds their difference uniformly
on the whole compact set. Completeness of \(\mathbb C\) gives a
uniform limit. The displayed modulus makes the limit continuous. Limits
on overlapping compact neighborhoods agree by pointwise uniqueness,
giving a single continuous \(f\) on \(U\).

This reasoning also proves a useful moving-point assertion: if the
rounded \(G_N\) converge locally uniformly to a continuous \(g\),
and grid points \(y_N\to y\in U\), then \(G_N(y_N)\to g(y)\).
Indeed, eventually the points lie in one compact neighborhood; use
\(Q_Ny_N=y_N\), uniform convergence there, and continuity of \(g\).
The same statement holds after any indices \(N_k\to\infty\).

\emph{Cauchy convergence of the next differences.}
Suppose now the rounded \(G_N=F_{N,\alpha}\) converge locally uniformly
to \(g\). Fix \(x\in U\) and a coordinate \(j\). Take a small
\(h>0\) such that a compact neighborhood of the segment
\([x,x+he_j]\) lies in \(U\). Put
\(h_N=a_N\lceil h/a_N\rceil\), \(x_N=Q_Nx\), and
\(y_N=x_N+h_Ne_j\). Then \(h\le h_N<h+a_N\), so the moving-point
assertion implies convergence of
\[
 V_N=\frac{G_N(y_N)-G_N(x_N)}{h_N}
 \longrightarrow\frac{g(x+he_j)-g(x)}h.
\]
The Taylor estimate gives
\[
 |\Delta_{N,j}G_N(x_N)-V_N|\le B_2h_N.
\]
Here \(B_2\) can be fixed on one compact neighborhood for every
sufficiently small \(h\). Given \(\epsilon>0\), first choose \(h\)
so that \(4B_2h<\epsilon/2\), then choose \(N\) large so that
\(h_N\le2h\) and \(V_N\) is Cauchy with error less than
\(\epsilon/2\). The triangle inequality proves that
\(\Delta_{N,j}G_N(Q_Nx)\) is Cauchy. Bounds for its first differences
give asymptotic equicontinuity by the previous argument. It therefore
converges locally uniformly to a continuous function \(g_j\).

\emph{Identification of derivatives and induction.}
Apply \eqref{ex-grid-taylor} with \(Q_Nx,Q_Ny\), for fixed nearby
\(x,y\), and pass to the limit using the just-established convergence
of all first differences. This gives
\[
 |g(y)-g(x)-\sum_j(y_j-x_j)g_j(x)|\le B_2|y-x|_1^2.
\]
The bound proves Fr\'echet differentiability, with
\(Dg(x)[u]=\sum_ju_jg_j(x)\). Since the \(g_j\) are continuous,
\(g\) is \(C^1\). Repeat this argument for every finite list
\(\alpha\). Induction gives locally uniform limits for all iterated
differences and identifies them with the corresponding derivatives of
\(f\). Their continuity proves \(f\in C^k\) for every finite \(k\),
hence smoothness. The case \(d=0\) is simply convergence of one scalar
sequence; there are no derivative steps. This completes the proof.
\end{proof}

\begin{corollary}[Other fixed difference stencils]\label{ex-other-stencils}
The same limits hold for any fixed finite composition of forward,
backward or centered coordinate differences, provided every difference
is divided by its full physical interval length.
\end{corollary}
\begin{proof}
For fixed integers \(b,c\ge0\) with \(b+c>0\), telescope:
\[
 \frac{G_N(x+ca_Ne_j)-G_N(x-ba_Ne_j)}{(b+c)a_N}
 =\frac1{b+c}\sum_{r=0}^{b+c-1}
       \Delta_{N,j}G_N(x+(r-b)a_Ne_j).
\]
Each of the finitely many sample points approaches the same point, so
the moving-point conclusion in Theorem~\ref{ex-smooth-grid-limit} gives
the same derivative for each summand. Translation averages commute with
coordinate differences. A finite composition is therefore a finite
average of shifted iterated unit forward differences, and has exactly
the stated limiting derivative. Uniform convergence on compacts follows
by using one slightly enlarged compact neighborhood for all the finitely
many shifts. No assertion is made when the stencil sizes vary with \(N\).
\end{proof}

The theorem is applied separately for each fixed period, label tuple and
finite derivative request. Its derivative bounds are supplied by the
full-precision lattice estimates, and its dense Cauchy input by the
renormalization-group comparison of different cutoffs. Neither input
follows from a subsequential limit alone. Once they are supplied, the
fixed-period smooth kernel and all its locally uniform difference limits
follow without further mesh extraction. The fixed-period smeared limits
then follow by the cell-integral dominated-convergence argument of
Proposition~\ref{prop:ex-limits-extraction}, applied to each fixed stencil.
